\documentclass[11pt]{article}
\usepackage[a4paper,margin=1in]{geometry}
\usepackage{amsmath,amssymb,amsthm,mathtools,bm}
\usepackage{hyperref}

\title{Integrability Assumption for Swapping Integral and Finite Sums\\in the QAOA--SK Expansion}
\author{}
\date{}

\begin{document}
	\maketitle
	
	\section*{Context}
	In the current Lean development, the quantity of interest is
	\begin{equation}
		\mathbb{E}_{J}\!\left[\langle s|U_C^\dagger U_B^\dagger \frac{C}{n+1} U_B U_C|s\rangle\right].
	\end{equation}
	After basis expansion and simplification of the two $|s\rangle$ amplitudes, we have
	\begin{equation}
		\mathbb{E}_{J}[\cdots]
		=
		\left(\frac12\right)^{n+1}
		\int
		\sum_{x,y,z}
		F_{x,y,z}(J)\,d\mu(J),
		\label{eq:before-swap}
	\end{equation}
	where (for fixed $x,y,z$)
	\begin{equation}
		F_{x,y,z}(J)
		=
		U_C^{\dagger\,\mathrm{phase}}(x;J)
		\,\overline{K_\beta(z,x)}
		\,\frac{C_J(z)}{n+1}
		\,K_\beta(z,y)
		\,U_C^{\mathrm{phase}}(y;J).
	\end{equation}
	
	\section*{What the assumption says}
	The theorem \texttt{expectedQAOACostVEV\_swap\_integral\_sum} introduces
	\begin{equation}
		\forall x,y,z,\quad F_{x,y,z} \in L^1(\mu),
	\end{equation}
	i.e. each indexed summand is Bochner-integrable (here codomain is $\mathbb{C}$, so this is the usual complex integrability condition):
	\begin{equation}
		\int \lvert F_{x,y,z}(J)\rvert\,d\mu(J) < \infty,
	\end{equation}
	and measurability.
	
	In Lean, this is the hypothesis
	\begin{equation}
		\texttt{hInt :}\; \forall x\,y\,z,\; \texttt{MeasureTheory.Integrable (fun J => ... ) $\mu$}.
	\end{equation}
	
	\section*{Why this assumption is needed}
	To exchange integral and finite sums, the proof repeatedly applies
	\begin{equation}
		\texttt{MeasureTheory.integral\_finset\_sum}.
	\end{equation}
	That lemma requires integrability of each function being summed. Since we have a triple finite sum, we apply the lemma three times (for $x$, then $y$, then $z$). The assumption \texttt{hInt} supplies exactly these hypotheses.
	
	\section*{Result after swapping}
	Under \texttt{hInt}, equation \eqref{eq:before-swap} becomes
	\begin{equation}
		\mathbb{E}_{J}[\cdots]
		=
		\left(\frac12\right)^{n+1}
		\sum_{x,y,z}
		\int F_{x,y,z}(J)\,d\mu(J).
		\label{eq:after-swap}
	\end{equation}
	This is precisely the mathematical content of the Lean theorem
	\texttt{expectedQAOACostVEV\_swap\_integral\_sum}.
	
	\section*{No extra physics assumptions}
	This step does \emph{not} use SK-specific Gaussian identities yet. It is a measure-theoretic rearrangement step:
	\begin{itemize}
		\item algebraic structure from earlier expansions, and
		\item integrability to justify changing order of finite sums and integration.
	\end{itemize}
	The SK/Gaussian calculations enter in later steps when evaluating each integral term in \eqref{eq:after-swap} analytically.
	
\end{document}
